% Generado por semana10_tablas_paper.py a partir de results/. NO editar a mano.

\begin{table*}[t]\centering\footnotesize\setlength{\tabcolsep}{4pt}
\caption{O1/H1. AP por bloque en prueba bajo el Esquema B (media de tres semillas por bloque; nivel de azar = prevalencia media de los bloques, 0.031). $\Delta$: diferencia media pareada frente a G0 con IC 95\% bootstrap; $p_{\mathrm{H}}$: Wilcoxon unilateral con Holm, contra $\delta=0$ y contra $\delta=0.05$; Gana: bloques en que el detector supera a G0.}
\label{tab:h1}
\begin{tabular}{@{}lccccccccc@{}}\toprule
Detector & $n$ & Media & DE & Mediana & IC 95\% & $\Delta$ [IC 95\%] & $p_{\mathrm{H}}$ ($\delta{=}0$) & $p_{\mathrm{H}}$ ($\delta{=}0.05$) & Gana \\\midrule
G0 m\'ovil & 291 & 0.345 & 0.080 & 0.348 & 0.336--0.354 & -- & -- & -- & -- \\
G1 LSTM & 291 & 0.997 & 0.008 & 0.999 & 0.996--0.998 & +0.652 [+0.642, +0.661] & $<\!10^{-10}$ & $<\!10^{-10}$ & 291/291 \\
G2 GRU & 291 & 0.996 & 0.007 & 0.999 & 0.996--0.997 & +0.651 [+0.642, +0.660] & $<\!10^{-10}$ & $<\!10^{-10}$ & 291/291 \\
G3 CNN-1D & 291 & 0.995 & 0.010 & 0.998 & 0.994--0.996 & +0.650 [+0.641, +0.659] & $<\!10^{-10}$ & $<\!10^{-10}$ & 291/291 \\
G4 AE & 291 & 0.449 & 0.168 & 0.459 & 0.430--0.467 & +0.103 [+0.082, +0.123] & $<\!10^{-10}$ & $3.7\!\times\!10^{-7}$ & 212/291 \\
G5 Transf. & 291 & 0.997 & 0.007 & 1.000 & 0.996--0.998 & +0.652 [+0.643, +0.661] & $<\!10^{-10}$ & $<\!10^{-10}$ & 291/291 \\
G6 CNN-LSTM-GRU & 291 & 0.997 & 0.006 & 0.999 & 0.996--0.998 & +0.652 [+0.642, +0.661] & $<\!10^{-10}$ & $<\!10^{-10}$ & 291/291 \\
\bottomrule\end{tabular}\end{table*}

\begin{table*}[t]\centering\scriptsize\setlength{\tabcolsep}{2.6pt}
\caption{M\'etricas en prueba. Esquema A: solo descriptivo (comparte familia de reglas con G0); 84 bloques con ambas clases; azar = prevalencia media de esos bloques (0.0034; la global es 0.0010). Esquema B: 291 bloques, azar 0.031; AP con DE e IC 95\% en la Tabla~\ref{tab:h1}. F1 y FPR dependen de umbrales elegidos en el bloque de umbral, que forma parte del entrenamiento final y puede favorecerlos; el AP no. R90: umbral que alcanz\'o recall 0.90 en ese bloque; F1: umbral de F1 m\'aximo.}
\label{tab:secund}
\begin{tabular}{@{}l|ccccccccc|ccccc@{}}\toprule
& \multicolumn{9}{c|}{Esquema A ($n=84$)} & \multicolumn{5}{c}{Esquema B ($n=291$)} \\
Detector & AP & DE & IC 95\% & Med. & $\Delta$ & F1 & FPR F1 & Rec.\ R90 & FPR R90 & ROC & F1 & FPR F1 & Rec.\ R90 & FPR R90 \\\midrule
G0 m\'ovil & 0.084 & 0.195 & 0.046--0.129 & 0.014 & -- & 0.021 & 0.031 & 0.976 & 0.619 & 0.792 & 0.323 & 0.027 & 0.899 & 0.546 \\
G1 LSTM & 0.416 & 0.296 & 0.363--0.477 & 0.457 & +0.333 & 0.213 & $7\!\times\!10^{-4}$ & 0.932 & 0.065 & 1.000 & 0.988 & $3\!\times\!10^{-4}$ & 0.902 & $2\!\times\!10^{-5}$ \\
G2 GRU & 0.407 & 0.286 & 0.343--0.467 & 0.425 & +0.323 & 0.233 & $7\!\times\!10^{-4}$ & 0.929 & 0.067 & 1.000 & 0.987 & $3\!\times\!10^{-4}$ & 0.900 & $2\!\times\!10^{-5}$ \\
G3 CNN-1D & 0.348 & 0.266 & 0.293--0.400 & 0.355 & +0.265 & 0.192 & $10\!\times\!10^{-4}$ & 0.888 & 0.148 & 0.999 & 0.985 & $3\!\times\!10^{-4}$ & 0.898 & $3\!\times\!10^{-5}$ \\
G4 AE & 0.038 & 0.084 & 0.023--0.057 & 0.010 & -0.046 & 0.035 & 0.123 & 0.982 & 0.826 & 0.922 & 0.421 & 0.045 & 0.906 & 0.196 \\
G5 Transf. & 0.371 & 0.283 & 0.312--0.431 & 0.356 & +0.288 & 0.204 & $10\!\times\!10^{-4}$ & 0.953 & 0.141 & 1.000 & 0.986 & $3\!\times\!10^{-4}$ & 0.907 & $3\!\times\!10^{-5}$ \\
G6 CNN-LSTM-GRU & 0.399 & 0.315 & 0.332--0.465 & 0.415 & +0.315 & 0.218 & $8\!\times\!10^{-4}$ & 0.925 & 0.210 & 1.000 & 0.988 & $3\!\times\!10^{-4}$ & 0.903 & $3\!\times\!10^{-5}$ \\
\bottomrule\end{tabular}\end{table*}

\begin{table*}[t]\centering\footnotesize\setlength{\tabcolsep}{4pt}
\caption{AP medio por bloque en prueba (Esquema B): por semilla ($n=291$) y por tipo de anomal\'ia, con la DE entre par\'entesis ($n=291$, $291$ y $291$ bloques con instantes puntuales, contextuales y colectivos). Por tipo: positivos del tipo frente a todos los instantes normales; los positivos de los otros tipos se excluyen. La \'ultima fila es el nivel de azar de cada tipo.}
\label{tab:tipo}
\begin{tabular}{@{}lcccccc@{}}\toprule
Detector & Sem.\ 42 & Sem.\ 43 & Sem.\ 44 & Puntual & Contextual & Colectiva \\\midrule
G0 m\'ovil & -- & -- & -- & 0.905 (0.103) & 0.399 (0.197) & 0.131 (0.057) \\
G1 LSTM & 0.997 & 0.997 & 0.997 & 0.999 (0.003) & 0.947 (0.079) & 0.999 (0.009) \\
G2 GRU & 0.996 & 0.997 & 0.996 & 0.999 (0.003) & 0.936 (0.086) & 0.999 (0.007) \\
G3 CNN-1D & 0.995 & 0.996 & 0.994 & 0.995 (0.021) & 0.931 (0.091) & 0.997 (0.011) \\
G4 AE & 0.443 & 0.450 & 0.452 & 0.733 (0.275) & 0.016 (0.019) & 0.340 (0.179) \\
G5 Transf. & 0.997 & 0.998 & 0.997 & 0.998 (0.008) & 0.955 (0.070) & 0.998 (0.008) \\
G6 CNN-LSTM-GRU & 0.997 & 0.997 & 0.997 & 0.998 (0.014) & 0.941 (0.090) & 0.999 (0.005) \\
\midrule Azar & -- & -- & -- & 0.005 & 0.003 & 0.023 \\
\bottomrule\end{tabular}\end{table*}

\begin{table}[t]\centering\footnotesize\setlength{\tabcolsep}{3pt}
\caption{Ablaci\'on de canales (Esquema B, prueba, una semilla): AP medio por bloque con ocupaci\'on y velocidad en cero (solo flujo y variables temporales), frente a la corrida completa y a G0.}
\label{tab:abl}
\begin{tabular}{@{}lccccccc@{}}\toprule
Detector & Completa & Ablada & $\Delta$ & $n$ & Punt. & Context. & Colect. \\\midrule
G1 LSTM & 0.997 & 0.879 & -0.118 & 291 & 0.935 & 0.664 & 0.843 \\
G2 GRU & 0.996 & 0.843 & -0.153 & 291 & 0.942 & 0.639 & 0.795 \\
G3 CNN-1D & 0.995 & 0.939 & -0.055 & 291 & 0.931 & 0.626 & 0.936 \\
G5 Transf. & 0.997 & 0.948 & -0.049 & 291 & 0.933 & 0.696 & 0.938 \\
G6 CNN-LSTM-GRU & 0.997 & 0.955 & -0.042 & 291 & 0.941 & 0.684 & 0.949 \\
G0 m\'ovil (solo flujo) & -- & 0.345 & -- & 291 & 0.905 & 0.399 & 0.131 \\
\bottomrule\end{tabular}\end{table}
